% 本套各册使用同一字号与页式；不对整页做缩放。
\usepackage[left=14mm,right=14mm,top=19mm,bottom=17mm,headheight=16pt,headsep=5mm,footskip=9mm]{geometry}
\usepackage{amsmath,amssymb,array,graphicx,xcolor,tikz,adjustbox,hyperref,needspace,booktabs,tabularx,enumitem,fancyhdr,etoolbox,multicol}
\usetikzlibrary{calc,angles,quotes,arrows.meta,patterns,decorations.pathreplacing,positioning}
\xeCJKDeclareCharClass{CJK}{"2460 -> "2473}
\definecolor{answerred}{RGB}{196,20,38}
\definecolor{bookred}{gray}{0}
\definecolor{bookpink}{gray}{1}
\definecolor{bookblue}{gray}{0}
\setlength{\parindent}{0pt}
\setlength{\parskip}{2pt}
\setlength{\emergencystretch}{3em}
\setlength{\abovedisplayskip}{3pt}\setlength{\belowdisplayskip}{3pt}
\setlength{\abovedisplayshortskip}{2pt}\setlength{\belowdisplayshortskip}{2pt}
\setlength{\jot}{1pt}
\setlength{\columnsep}{7mm}\setlength{\columnseprule}{.25pt}
\setlist{nosep,leftmargin=1.8em,topsep=2pt,itemsep=2pt}
\hypersetup{unicode=true,hidelinks,pdfcreator={XeLaTeX / ctexbook},bookmarksnumbered=false}
\newcommand{\SourceNumber}{1}\newcommand{\SourceTopic}{}
\pagestyle{fancy}\fancyhf{}
\fancyhead[L]{\normalsize\SourceTopic}
\fancyhead[C]{\normalsize 原书第\SourceNumber 页}
\fancyfoot[C]{\normalsize\thepage}
\renewcommand{\headrulewidth}{.35pt}
\fancypagestyle{plain}{\fancyhf{}\fancyhead[L]{\normalsize\SourceTopic}\fancyhead[C]{\normalsize 原书第\SourceNumber 页}\fancyfoot[C]{\normalsize\thepage}\renewcommand{\headrulewidth}{.35pt}}
\newcommand{\SourcePage}[2]{\clearpage\gdef\SourceNumber{#1}\gdef\SourceTopic{#2}\phantomsection\label{src:#1}\pdfbookmark[0]{原书第#1页}{source-#1}}
\newcommand{\BookPage}[2]{}\newcommand{\FirstBookPage}[2]{}\newcommand{\bookpage}[1]{}
\newcommand{\QHeading}[1]{\par\addvspace{6pt}{\heiti\normalsize\color{black}#1}\par\nobreak\vspace{2pt}}
\ctexset{subsection={format=\normalsize\heiti,beforeskip=6pt,afterskip=2pt}}
\newcommand{\Problem}[1]{\par\noindent{\color{black}#1}\par\nobreak}
\newcommand{\PlainItem}[1]{\par\noindent{\color{black}#1}\par}
\newcommand{\Answer}[1]{\par\begingroup\color{answerred}#1\par\endgroup}
\newcommand{\Formula}[1]{\par\noindent$\displaystyle#1$\par}
\newcommand{\Equation}[1]{\Formula{#1}}
\newcommand{\CalcBlock}[3][]{\par\noindent\begin{minipage}[t]{\linewidth}\color{black}（#2）\ifstrempty{#1}{}{\quad\textcolor{answerred}{#1}}\par\nobreak\Formula{#3}\end{minipage}\par}
\newcommand{\MethodTitle}[1]{\par{\color{answerred}\heiti #1}\par\nobreak}
\newcommand{\MethodBlock}[3]{\par\noindent\begin{minipage}{\linewidth}\color{answerred}{\heiti #1}\par #2\par\Formula{#3}\end{minipage}\par}
\newenvironment{oralgrid}{\par\begingroup\setlength{\parskip}{2pt}}{\par\endgroup}
\newcommand{\OralItem}[3]{\noindent\begin{minipage}[t]{.245\linewidth}\strut$#2=\textcolor{answerred}{#3}$\strut\end{minipage}\allowbreak\ignorespaces}
\newcommand{\DataTable}[2]{\par\begingroup\setlength{\tabcolsep}{3pt}\renewcommand{\arraystretch}{1.15}\noindent\begin{tabularx}{\linewidth}{|*{#1}{>{\centering\arraybackslash}X|}}\hline #2\end{tabularx}\par\endgroup}
\newcommand{\q}[5]{\node[anchor=north west,inner sep=0pt,outer sep=0pt,align=left,text width=\dimexpr#3bp*2/3\relax,font=\normalsize,text=black] at ({#1/1.5},{#2/1.5}) {#5};}
\newcommand{\ans}[5]{\node[anchor=north west,inner sep=0pt,outer sep=0pt,align=left,text width=#3bp,font=\normalsize,text=answerred] at (#1,#2) {#5};}
\newcommand{\hdhead}[2]{}\newcommand{\hdsection}[3]{}
\tikzset{every picture/.style={line cap=round,line join=round,font=\normalsize}, hdheight/.style={draw=answerred,dashed,line width=.85pt},hdgrid/.style={draw=black,line width=.23pt},hdoutline/.style={draw=black,line width=.65pt},outline/.style={draw=black,line width=.65pt},guide/.style={draw=black,line width=.28pt},axis/.style={draw=black,dash pattern=on 3pt off 2pt,line width=.5pt},thinrule/.style={draw=black,line width=.4pt},measure/.style={draw=black,line width=.4pt},number/.style={font=\normalsize,inner sep=1pt}}
\newcommand{\keyheading}[1]{\QHeading{#1}}
\newcommand{\keygrid}[2]{\par\begingroup\renewcommand{\arraystretch}{1.2}\noindent\begin{tabularx}{\linewidth}{@{}*{#1}{>{\centering\arraybackslash}X}@{}}#2\end{tabularx}\par\endgroup}
\newcommand{\WorkSpace}[1]{\par\vspace{#1mm}}
\newenvironment{QuestionGrid}[1]{\par\noindent\begin{tabular}{@{}*{#1}{p{\dimexpr(\linewidth-12pt*#1+12pt)/#1\relax}}@{}}}{\end{tabular}\par}
\newcommand{\MathCell}[3]{\begin{minipage}[t]{\linewidth}\strut #1$\displaystyle #2$\par\vspace{#3mm}\end{minipage}}
\newcommand{\QuestionFigure}[1]{\par\noindent\input{#1}\par}
\newsavebox{\SourceBox}
\newcommand{\OnePageInput}[1]{\setbox\SourceBox=\vbox{\hsize=\textwidth\linewidth=\textwidth\columnwidth=\textwidth\normalsize\input{#1}\par}\typeout{SOURCEHEIGHT \SourceNumber: \the\ht\SourceBox / \the\textheight}\nointerlineskip\box\SourceBox}
\renewcommand{\Needspace}[1]{}\renewcommand{\needspace}[1]{}
\raggedbottom\clubpenalty=10000\widowpenalty=10000

% 原书第77页的独立矢量题图。
\input{../p77_tikz.tex}
